\section{Phase 4: nonlinear-plant escalation}
\label{sec:phase4}
Phase~4 replaces the linear plant, not the compact identified model class. It tests
whether the structural conclusion survives realistic plant--model mismatch.

\subsection{Experiment 4A: nonlinear tire-model validation}
\paragraph{Objective.}
Implement and validate saturating tire forces
\begin{equation}
 F_{y,j}=C_j\alpha_{\mathrm{sat},j}
 \tanh\left(\frac{\alpha_j}{\alpha_{\mathrm{sat},j}}\right),
 \qquad j\in\{f,r\}.
\end{equation}

\paragraph{Predeclared expectation.}
The nonlinear plant agrees with the linear model at small slip and departs smoothly
as tire saturation becomes relevant.

\paragraph{Setup and planned evidence.}
Check the force--slip curves, small-signal slope, force limits, integration
accuracy, sideslip, yaw rate, and lateral acceleration. Preserve the Phase-0
structural families initially so that the plant change is isolated.

\paragraph{Decision gate.}
The model must have correct small-slip behavior, smooth saturation, plausible force
limits, and numerically converged simulations. Otherwise Phase~4 is blocked.

\paragraph{Results, interpretation, and decision.}
The nonlinear plant retains the Phase-0 state, input, structural family
parameters, and sign conventions. A common dry-road friction coefficient
$\mu=0.9$ is used. Static axle loads are
\begin{equation}
 F_{z,f}=mg\frac{l_r}{l_f+l_r},\qquad
 F_{z,r}=mg\frac{l_f}{l_f+l_r},
\end{equation}
and the saturation angles are selected as
\begin{equation}
 \alpha_{\mathrm{sat},j}=\frac{\mu F_{z,j}}{C_j}.
\end{equation}
Consequently, each tanh law has the original cornering stiffness at zero slip
and approaches the axle friction limit $\mu F_{z,j}$ smoothly. The resulting
family-center scales and force limits are listed in \cref{tab:experiment4a-tires}.

\begin{table}[htbp]
\centering
\caption{Experiment 4A tanh-tire parameters for the three Phase-0 family
centers at $\mu=0.9$.}
\label{tab:experiment4a-tires}
\begin{tabular}{lrrrr}
\toprule
Family & $\alpha_{\mathrm{sat},f}$ [deg] & $\alpha_{\mathrm{sat},r}$ [deg]
& $F_{y,f}^{\max}$ [kN] & $F_{y,r}^{\max}$ [kN]\\
\midrule
Nominal & 5.420 & 4.065 & 7.568 & 5.676\\
Heavy & 5.936 & 4.675 & 8.703 & 6.527\\
Handling & 5.721 & 3.511 & 7.189 & 5.392\\
\bottomrule
\end{tabular}
\end{table}

Central finite differences were used to linearize the nonlinear right-hand side
at zero sideslip, yaw rate, and steering for all family centers at
$v_x\in\{12,20,28\}\,\mathrm{m\,s^{-1}}$. The maximum relative Jacobian error
against the analytical Phase-0 matrices is $9.74\times10^{-13}$. Numerical force
curves are odd, monotone, have the prescribed zero-slip slope, and remain within
their friction limits.

Open-loop simulations use a smooth steering doublet at the same three speeds.
Three diagnostic amplitudes expose the transition from the linear limit to tire
saturation. Their worst results over all family--speed combinations are reported
in \cref{tab:experiment4a-responses}. These amplitudes validate the plant only;
Experiment 4B will independently select the final near-linear, moderate, and
strong-but-feasible maneuver levels.

\begin{table}[htbp]
\centering
\caption{Worst Experiment 4A nonlinear response diagnostics over the three
family centers and speeds $12$, $20$, and $28\,\mathrm{m\,s^{-1}}$. The model
difference is trajectory RMSE relative to the corresponding linear response.}
\label{tab:experiment4a-responses}
\begin{tabular}{lrrrr}
\toprule
Input & Amplitude [deg] & Model difference [\%] & Force utilization & $|a_y|_{\max}$ [$\mathrm{m\,s^{-2}}$]\\
\midrule
Small & 0.1 & 0.0084 & 0.032 & 0.274\\
Moderate & 2.0 & 3.705 & 0.595 & 5.121\\
Strong & 6.0 & 173.7 & 1.000 & 8.826\\
\bottomrule
\end{tabular}
\end{table}

The small input produces virtually identical linear and nonlinear trajectories.
The moderate input introduces a visible but controlled departure. The strongest
case drives at least one axle close to its friction limit and therefore separates
sharply from the unbounded linear-tire prediction, especially at high speed. All
27 nonlinear family--speed--amplitude simulations remain finite. These behaviors,
together with the tire curves, are summarized in
\cref{fig:experiment4a-validation}.

\begin{figure}[htbp]
 \centering
 \includegraphics[width=0.95\linewidth]{experiment_4a_nonlinear_validation.pdf}
 \caption{Experiment 4A nonlinear-plant validation. The tanh tire law recovers
 the prescribed small-slip slope and saturates smoothly. The nonlinear vehicle
 matches the linear model under weak excitation and departs under strong
 excitation. The RK4 trajectory error decreases at the expected fourth-order
 rate.}
 \label{fig:experiment4a-validation}
\end{figure}

Integration was audited for the nominal $20\,\mathrm{m\,s^{-1}}$, $6^\circ$
case using RK4 steps from $40$ to $5\,\mathrm{ms}$ against a $1.25\,\mathrm{ms}$
reference. The relative trajectory error decreases by approximately a factor of
16 whenever the step is halved. At the retained $10\,\mathrm{ms}$ step it is
$7.10\times10^{-7}\%$, well below the predeclared $0.1\%$ gate.

Experiment 4A therefore \textbf{passes}. The nonlinear equations recover the
validated linear plant locally, provide bounded and physically interpretable tire
forces, and are numerically converged at the existing sample time. Experiment 4B
may proceed without changing the Phase-0 family definitions.
\subsection{Experiment 4B: maneuver-severity calibration}
\paragraph{Objective.}
Define reproducible near-linear, moderately nonlinear, and strongly nonlinear but
feasible operating levels.

\paragraph{Setup and planned evidence.}
Scale the lateral reference by
$\gamma_m\in\{0.5,1.0,1.5\}$. Report peak slip, lateral acceleration, tire-force
utilization, steering limits, and the fraction of time outside a predeclared
near-linear tire region. Adjust the absolute base maneuver once, globally, if the
three intended regimes are not obtained.

\paragraph{Decision gate.}
Severity levels must represent distinct regimes without making the strongest case
physically meaningless or infeasible for every linear controller.

\paragraph{Results, interpretation, and decision.}
Pending.

\subsection{Experiment 4C: oracle models under nonlinear mismatch}
\paragraph{Objective.}
Repeat M1--M5 using nonlinear-plant trajectories while retaining linear LTI/LPV
identified models.

\paragraph{Predeclared expectation.}
Absolute errors grow with maneuver severity, but the complementary benefits of
scheduling and specialization remain under moderate nonlinear excitation.

\paragraph{Setup and decision gate.}
Evaluate one-step and free-run error versus speed, family, and severity. The M4
advantage must survive at least the moderate regime and unseen speed--maneuver
combinations. If all linear LPV models fail, revise the representation before FL.

\paragraph{Results, interpretation, and decision.}
Pending.

\subsection{Experiment 4D: controller transfer and redesign}
\paragraph{Objective.}
Separate robustness of the Phase-3 controller from the achievable performance of
controllers redesigned using LPV approximations identified on nonlinear data.

\paragraph{Setup.}
The transfer test applies Phase-3 controllers unchanged. The redesign test refits
the compact models on nonlinear trajectories and synthesizes controllers with the
same common weights. The redesign test is primary; transfer is diagnostic.

\paragraph{Decision gate.}
The qualitative M4 benefit should survive moderate excitation. If it survives only
near the linear region, the claim and operating envelope are restricted
explicitly. Frozen-grid checks are necessary but are not presented as a formal LPV
stability certificate.

\paragraph{Results, interpretation, and decision.}
Pending.
